\documentclass[12pt,a4paper,UTF8]{ctexart}
\usepackage[a4paper,top=22mm,bottom=22mm,left=22mm,right=22mm,headheight=15pt]{geometry}
\newcommand{\Needspace}[1]{\par\begingroup\dimen0=#1\relax\dimen2=\pagegoal\advance\dimen2 by-\pagetotal\ifdim\dimen0>\dimen2\newpage\fi\endgroup}
\newcommand{\needspace}[1]{\par\begingroup\skip0=\lastskip\vskip-\skip0\dimen0=#1\relax\advance\dimen0 by\skip0\vskip\dimen0\penalty-100\vskip-\dimen0\vskip\skip0\penalty9999\endgroup}
\usepackage{amsmath,amssymb,mathtools,bm}
\usepackage{graphicx,xcolor,booktabs,longtable,array,enumitem,fancyhdr,caption,float,microtype,setspace}
\usepackage[hidelinks,unicode]{hyperref}
\usepackage[most]{tcolorbox}
\definecolor{ContestBlue}{HTML}{17365D}
\hypersetup{pdftitle={第072套},pdfauthor={}}
\setstretch{1.24}
\setlength{\parindent}{2em}
\setlength{\parskip}{0.35em}
\setlength{\emergencystretch}{2em}
\setlist{itemsep=0.35em,topsep=0.4em}
\captionsetup{font=small,labelfont=bf,justification=centering}
\setcounter{secnumdepth}{0}
\pagestyle{fancy}
\fancyhf{}
\fancyhead[L]{\small 第072套}
\fancyfoot[C]{\thepage}
\renewcommand{\headrulewidth}{0.4pt}
\ctexset{
  section={format=\Large\bfseries\color{ContestBlue},beforeskip=1.2em,afterskip=0.7em},
  subsection={format=\large\bfseries\color{ContestBlue},beforeskip=1.0em,afterskip=0.5em},
  subsubsection={format=\normalsize\bfseries,beforeskip=0.8em,afterskip=0.35em}
}
\usepackage{calc}
\newcounter{none}
\let\originalsubsection\subsection
\renewcommand{\subsection}{\needspace{6\baselineskip}\originalsubsection}
\makeatletter
\newcommand{\tightlist}{\setlength{\itemsep}{0pt}\setlength{\parskip}{0pt}}
\newcommand{\maxwidth}{\ifdim\Gin@nat@width>\linewidth\linewidth\else\Gin@nat@width\fi}
\newcommand{\maxheight}{\ifdim\Gin@nat@height>0.72\textheight 0.72\textheight\else\Gin@nat@height\fi}
\newcommand{\pandocbounded}[1]{\begingroup\setkeys{Gin}{width=\maxwidth,height=\maxheight,keepaspectratio}#1\endgroup}
\makeatother
\begin{document}
\begin{titlepage}
  \centering
  \vspace*{58mm}
  {\Huge\bfseries\color{ContestBlue} 第072套\par}
  \vfill
  \vspace*{28mm}
\end{titlepage}
\section{第 1 题}\label{ux7b2c-1-ux9898}

\begin{enumerate}
\def\labelenumi{\arabic{enumi}.}
\tightlist
\item
  光滑长斜面倾角 \(30^\circ\)，同一直线上放置 \(10\)
  个完全相同的滑块。解除约束时分别给予任意大小、可沿斜面向上或向下的初速度。所有碰撞完全弹性，求滑块之间最多碰撞次数。
\item
  光滑水平面左端有墙，质量 \(M=10^5\,\mathrm{kg}\) 的大滑块在右，质量
  \(m=10^{-5}\,\mathrm{kg}\)
  的小滑块在墙与大滑块之间。小滑块初静止，大滑块以
  \(v_0=3\,\mathrm{m\,s^{-1}}\) 向左。所有碰撞完全弹性，求总碰撞次数。
\end{enumerate}

\begin{figure}[H]
\centering
\includegraphics[width=0.84\linewidth,height=\textheight,keepaspectratio,alt={图1-1　同质量斜面滑块与墙边双质量碰撞系统}]{../figures/question-01.pdf}
\caption*{图1-1　同质量斜面滑块与墙边双质量碰撞系统}
\end{figure}

\newpage

\section{第 2 题}\label{ux7b2c-2-ux9898}

两个相同的大长方形绝缘平板平行，相距 \(d\)。上板面电荷密度为
\(+\sigma\)，下板为 \(-\sigma\)。点 \(P\) 位于上板一侧长边中点的正上方
\(h\) 处，且

\[
d\ll h\ll \text{平板边长}.
\]

估算 \(P\) 点电场大小和方向。

\begin{figure}[H]
\centering
\includegraphics[width=0.8\linewidth,height=\textheight,keepaspectratio,alt={图2-1　相反带电平行板边缘上方的观察点}]{../figures/question-02.pdf}
\caption*{图2-1　相反带电平行板边缘上方的观察点}
\end{figure}

\newpage

\section{第 3 题}\label{ux7b2c-3-ux9898}

两根高 \(50\,\mathrm m\) 的电线杆之间悬挂一根均匀柔软电线，线长
\(80\,\mathrm m\)，两端等高，最低点离地 \(20\,\mathrm m\)。

\begin{enumerate}
\def\labelenumi{\arabic{enumi}.}
\tightlist
\item
  求两杆间距。
\item
  求电线最低点的曲率半径。
\end{enumerate}

并整理 \(\sinh x,\cosh x\) 的恒等式、反函数及导数。

\begin{figure}[H]
\centering
\includegraphics[width=0.533\linewidth,height=\textheight,keepaspectratio,alt={图3-1　等高悬点之间的均匀悬链线}]{../figures/question-03.pdf}
\caption*{图3-1　等高悬点之间的均匀悬链线}
\end{figure}

\newpage

\section{第 4 题}\label{ux7b2c-4-ux9898}

过椭圆内部固定点 \(P\) 作一条动直线，与椭圆交于 \(A,B\)。在 \(A,B\)
处分别作椭圆切线，交于 \(Q\)。证明 \(Q\) 的轨迹在一条固定直线上。

\begin{figure}[H]
\centering
\includegraphics[width=0.72\linewidth,height=\textheight,keepaspectratio,alt={图4-1　过椭圆内定点的动弦及两端切线交点}]{../figures/question-04.pdf}
\caption*{图4-1　过椭圆内定点的动弦及两端切线交点}
\end{figure}

\newpage

\section{第 5 题}\label{ux7b2c-5-ux9898}

三根导线分别沿正方体棱组成图示闭合回路，长度依次为 \(4a,6a,6a\)。已知图
a 回路自感为 \(L_1\)，图 b 回路自感为 \(L_2\)，求图 c 回路自感。

\begin{figure}[H]
\centering
\includegraphics[width=0.88\linewidth,height=\textheight,keepaspectratio,alt={图5-1　正方体棱上的四棱、六棱与三面组合回路}]{../figures/question-05.pdf}
\caption*{图5-1　正方体棱上的四棱、六棱与三面组合回路}
\end{figure}

\newpage

\section{第 6 题}\label{ux7b2c-6-ux9898}

\(1\,\mathrm{mol}\) 铜块质量由摩尔质量 \(63.5\,\mathrm{g\,mol^{-1}}\)
给出，外包绝热膜。铜初温 \(27^\circ\mathrm C\)，摩尔热容
\(24.44\,\mathrm{J\,mol^{-1}K^{-1}}\)。质量 \(m=30\,\mathrm g\)
的铅弹初温 \(30^\circ\mathrm C\)，固态比热
\(128\,\mathrm{J\,kg^{-1}K^{-1}}\)，液态比热
\(153\,\mathrm{J\,kg^{-1}K^{-1}}\)，熔点
\(327.46^\circ\mathrm C\)，熔化热
\(2.045\times10^4\,\mathrm{J\,kg^{-1}}\)。

\begin{enumerate}
\def\labelenumi{\arabic{enumi}.}
\tightlist
\item
  铜块固定，铅弹以 \(850\,\mathrm{m\,s^{-1}}\) 射入并停止，求终温。
\item
  铜块在光滑水平面上，求终态热力学熵变。
\item
  铜块固定，速度改为 \(750\,\mathrm{m\,s^{-1}}\)，铅弹带电
  \(q=4.0\times10^{-5}\,\mathrm C\)。铅密度
  \(11.34\,\mathrm{g\,cm^{-3}}\)，终态电容
  \(C_f=1.5\,\mathrm{pF}\)。求熵变。
\end{enumerate}

\begin{figure}[H]
\centering
\includegraphics[width=0.86\linewidth,height=\textheight,keepaspectratio,alt={图6-1　铅弹嵌入绝热铜块的三种能量边界条件}]{../figures/question-06.pdf}
\caption*{图6-1　铅弹嵌入绝热铜块的三种能量边界条件}
\end{figure}

\newpage

\section{第 7 题}\label{ux7b2c-7-ux9898}

折射率为 \(n\)、折射棱角 \(\theta\ll1\) 的菲涅尔双棱镜，距轴上单色狭缝
\(S\) 为 \(a\)，距观察屏为 \(b\)，波长 \(\lambda\)。

\begin{enumerate}
\def\labelenumi{\arabic{enumi}.}
\tightlist
\item
  求相邻亮纹间距。
\item
  狭缝垂直主轴上移 \(h\)，求零级条纹移动方向和距离。
\item
  再贴合一个镜面对称的相同双棱镜，形成双面双棱镜，重做前两问。
\end{enumerate}

\begin{figure}[H]
\centering
\includegraphics[width=0.86\linewidth,height=\textheight,keepaspectratio,alt={图7-1　单面与双面菲涅尔双棱镜的虚光源}]{../figures/question-07.pdf}
\caption*{图7-1　单面与双面菲涅尔双棱镜的虚光源}
\end{figure}

\newpage

\section{第 8 题}\label{ux7b2c-8-ux9898}

\begin{enumerate}
\def\labelenumi{\arabic{enumi}.}
\tightlist
\item
  估算一个氩气分子线度的上限。
\item
  边长为 \(d\)
  的极小立方盒中只有一个中子。用两种方法估算盒壁压强的下限，已知普朗克常量
  \(h\) 和中子质量 \(m\)。
\end{enumerate}

\begin{figure}[H]
\centering
\includegraphics[width=0.8\linewidth,height=\textheight,keepaspectratio,alt={图8-1　单分子体积上限与立方盒中受限中子}]{../figures/question-08.pdf}
\caption*{图8-1　单分子体积上限与立方盒中受限中子}
\end{figure}
\end{document}
